\subsection{Neighborhoods and Interior}

\begin{definition}[Neighborhood System]
   Let \( (X,\mathcal{T}) \) be a topological space and \( x \in X  \). A \textbf{neighborhood} of \( x  \) is a subset of \( X  \) such that there exists \( A \in \mathcal{T} \) such that \( x \in A \subseteq U  \). The family of all neighborhoods of \( x  \) is called the \textbf{neighborhood system} of \( x  \) and is denoted by \( \mathcal{N}(x) \); that is,  
   \[  \mathcal{N}(x) = \{  U \subseteq X : U \ \text{is a neighborhood of \( x  \)} \}  \]
\end{definition}

\begin{problem}
    Consider \( \R  \) with the Euclidean topology and \( x \in \R  \).
    \begin{enumerate}
        \item[(i)] Give an example of a neighborhood of \( x  \) that is open and one that is not open.
            \begin{solution}
                \( (x - \epsilon, x + \epsilon) \) is an open neighborhood of \( x  \) with respect to the Euclidean Topology  and on the other hand, \( [x- \epsilon, x + \epsilon] \) is not an open neighborhood of \( x  \).
            \end{solution}
        \item[(ii)] Give an example of a subset of \( \R  \) containing \( x  \) that is not a neighborhood of \( x  \).
            \begin{solution}
                We have \( [x, x + 1) \), \( [x, x+ 1) \), and \( \{ x  \}  \).
            \end{solution}
    \end{enumerate}
\end{problem}

\begin{prop}
    Let \( X  \) be a topological space and \( A \subseteq X  \). We have that \( A  \) is open if and only if for all \( x \in A  \), we have \( A \in \mathcal{N}(x) \).
\end{prop}
\begin{proof}
\textit{Left to the reader.} \( \Longrightarrow  \) \textit{This is trivial.}
\( \Longleftarrow \)  Use the definition and apply the union and the union of open sets is open.
\end{proof}

\begin{definition}[Interior Point and Interior of a Set]
 Let \( X  \) be a topological space, \( A \subseteq  X  \), and \( x \in X  \).
   \begin{itemize}
       \item We say that \( x  \) in an \textbf{interior point} of \( A  \) if \( A \in \mathcal{N}(x) \); that is, there exists an open set \( U  \) such that \( x \in U \subseteq A  \).
        \item The \textbf{interior} of \( A  \) is defined as
            \[  \mathring{A} = \text{int}(A) = \{  x \in X : x \ \text{is an interior point of \( A  \)} \}.  \]
   \end{itemize} 
\end{definition}

\begin{remark}
    Let \( (X,\mathcal{T}) \) be a topological space and \( A \subseteq X  \). Then, 
    \begin{itemize}
        \item \( \mathring{A} \subseteq A  \),
        \item If \( \mathcal{B} \) is a basis for \( \mathcal{T} \), then \( x \in \mathring{A} \) if and only if there exists \( B \in \mathcal{B} \) such that \( x \in B \subseteq  A  \).
    \end{itemize}
    This allows us to verify if a point in the interior just by using the basis of \( \mathcal{T} \) instead of checking every open set. This makes it easier to check these properties because we are working with fewer sets.
\end{remark}

\begin{prop}
    Let \( X  \) be a topological space and \( A \subseteq X  \). We have that \( A  \) is open if and only if \( A = \mathring{A} \).
\end{prop}

\begin{proof}
    By the observation, \( A = \mathring{A} \) if and only if \( A \subseteq \mathring{A} \). Suppose \( A  \) is open. Then, \( A  \) is open if and only if for all \( x \in A  \), \( A  \) is a neighborhood of \( x  \) by proposition 1. Since \( A \in \mathcal{N}(x) \), for all \( x \in A  \), \( x \in \mathring{A} \) by definition of interior point. Hence, \( A \subseteq \mathring{A} \).
\end{proof}

What is the relevance of all of this? Notice that we declared that a set being open as an axiom. This proposition gives us a more deeper insight into what it means to be open.  


\begin{prop}
    Let \( X  \) be a topological space and \( A \subseteq X  \). The interior of \( A  \) is the largest open set contained in \( A  \). Equivalently, 
    \[  \mathring{A} = \bigcup_{ U \ \text{is open}, U \subseteq A  }^{  } U. \]
\end{prop}
\begin{proof}
    We proceed via double inclusion to prove the result. Denote by \( \mathcal{A} = \{  U \subseteq A : \text{\( U  \) is open} \}  \) Let \( x \in \mathring{A}  \). By defintion there exists an open set \( U  \) such that \( x \in U \subseteq A  \). Therefore, \( U \in \mathcal{A}  \) and so \( x \in \bigcup_{ U \in \mathcal{A} }^{  } U  \). This shows that the interior \( \mathring{A} \subseteq \bigcup_{ U \in \mathcal{A} }^{  } U  \). 

    Now, let \( x \in \bigcup_{ U \in \mathcal{A} }^{  } U  \). Then for some open set \( U \subseteq A  \), \( x \in U  \). By definition, this tells us that \( x  \) is an interior point of \( A  \). Hence, \( x \in \mathring{A} \). Thus, \( \bigcup_{ U \in \mathcal{A} }^{  } U \subseteq \mathring{A} \).
\end{proof}


\begin{problem}
    For \( a < b  \) in \( \R  \), consider \( A = (a,b) \), \( B = [a,b) \), \( C = [a,b] \), and \( D = \Q  \). Find the interior of \( A,B,C,D  \) as subsets of \( \R  \) endowed with 
    \begin{enumerate}
        \item[(i)] The Euclidean Topology,
            \begin{solution}
                \( A  \) itself is the interior because \( A  \) is open. \( A  \) is the interior of \( B \), \( A  \) is the interior of \( C \), then \( \mathring{D} = \emptyset \). We see this by contradiction. Suppose \( \mathring{Q} \neq \emptyset \). Then there exists an \( x \in \mathring{\Q} \). Then there exists an \( r > 0  \) such that \( x \in (x - r , x + r) \subseteq \Q  \). But this cannot be the case since this open set contains irrational points that are clearly not in \( \Q  \). Thus, a contradiction.
            \end{solution}
        \item[(ii)] The lower limit topology,
            \begin{solution}
                \begin{itemize}
                    \item Since \( B \) is open, then \( \mathring{B} = B = [a,b) \).
                    \item Recall that the lower limit topology is finer than the Euclidean topology (again because it has more open sets than the Euclidean topology), since \( A  \) is open in the Euclidean topology, it is also open in the lower limit topology.
      Thus, \( \mathring{A} = A = (a,b) \).
                    \item We claim that \( \mathring{C} = [a,b)  \). 

                        Let \( x \in C \). We will consider two cases; that is, if \( a \leq x < b  \). We need to find a \( U  \) such that \( x \in U \subseteq [a,b) \). Indeed, define \( U  = [x,b) \). Since \( x \in U  \) and \( U  \) is open with respect to our topology, it follows that \( x \in \mathring{C} \). Suppose now that \( x = b  \). Then this is not an interior point because there is no open set that contains \( b  \) is contained in \( [a,b) \); that is, for all sets \( U \in \mathcal{T} \), \( x \in U \not\subseteq [a,b)  \); that is, no interval of the form \( [c,d)  \) with \( c \leq x < d  \) is not contained in \( C  \) and we are done. 
                    \item Note that the interior of \( \Q  \) is just empty again.
                \end{itemize}
            \end{solution}
        \item[(iii)] The cofinite topology.
            \begin{solution}
                \textit{Left to reader.}   
            \end{solution}
    \end{enumerate}
\end{problem}

\begin{remark}
    Sets whose complements are finite have empty interiors.
\end{remark}

\subsection{Closed Sets and Closure}

\begin{definition}[Closed]
    Let \( X  \) be a topological space. A subset \( C  \) of \( X  \) is \textbf{closed} if its complement \( X \setminus  C  \) is open.
\end{definition}

\begin{eg}
    \begin{enumerate}
        \item[(i)] The sets \( [a,\infty) , (- \infty , a ] , [a,b]  \) for \( a,b \in \R  \) with \( a < b  \) are closed in \( \R  \) equipped with the Euclidean topology.
    \begin{itemize}
        \item \( \R \setminus  [a,+\infty) \) is open and so \( [a,+ \infty) \) is closed.
        \item \( R \setminus  (-\infty , a ]  = (a,\infty) \) is open, and so \( (-\infty  a ]  \) is closed.
        \item  \( \R \setminus  [a,b] = (-\infty , a) \cup (b,+\infty ) \) is open, so \( [a,b) \) is closed.
    \end{itemize}
        \item[(ii)] The sets \( [a,b] , (a,b]  \) for \( a,b \in \R  \) with \( a < b  \) are NOT closed in \( \R  \) equipped with the Euclidean topology.
        \item[(iii)] In a metric space, singletons are closed. \textit{proof left to reader.} 
    \end{enumerate}
\end{eg}

\begin{problem}
    Which of the sets in the previous example (part i) and ii)) are closed in \( \R  \) equipped with the lower limit topology?
    \textbf{ans:} The first and the third example. 

    We can actually write the these examples in terms of a union of open sets in the lower limit topology; that is, 
    \[  [a,\infty)^{c} = (-\infty , a ) = \bigcup_{ n \in \N  }^{  }  [a-n, a) \]
    so \( [a,\infty)^{c} \) is open, and so \( [a,\infty)   \) is closed.

    Also, we see that the second example is closed. Note 
    \[  (-\infty , a]^{c} = (a,\infty ) = \bigcup_{ n \in \N  }^{  }  \Big[a + \frac{ 1 }{ n }  , a + n \Big).  \]

    Note also that 
    \[  [a,b)^{c} = \underbrace{(-\infty , a )}_{\text{open}} \cup \underbrace{[b,\infty ]}_{\text{open}} \]
    so \( [a,b)^{c} \) is open and thus \( [a,b)  \) is closed.

    However, we see that \( (a,b]^{c} = (-\infty , a] \cup (b,+\infty )   \) not open (\textbf{hint:} look at \( a  \)). 

\end{problem}

\begin{problem}
    Which of the following sets are closed in \( \R^{2} \) equipped with the Euclidean topology?
\end{problem}
\begin{enumerate}
    \item[(i)] \( \{ ({x}_{1}, {x}_{2}) \in \R^{2} : {x}_{2} \leq  1  \}  \)

        Just show that the complement is open.
    \item[(ii)] \( \{ 0  \}  \times \R  \)
        
        Note that \( B^{c} = \{ ({x}_{1}, {x}_{2}) \in \R^{2} : {x}_{1} \neq 0  \}  = \{  ({x}_{1} , {x}_{2}) \in \R^{2} : {x}_{1} > 0  \}  \cup \{ ({x}_{1} , {x}_{2}) \in \R^{2} : {x}_{1} < 0  \}  \) which are relatively easy to show that both sets are open. 

        Alternatively, since we already know that singletons and the metric space are closed sets, then the finite product of closed sets is closed.
\end{enumerate}

\begin{prop}
    Let \( X  \) be a topological space. We have that 
    \begin{enumerate}
        \item[(i)] \( \emptyset  \) and \( X  \) are closed;
        \item[(ii)] If \( {C}_{1} \) and \( {C}_{2} \) are closed, then \( {C}_{1} \cup {C}_{2} \) is closed;
        \item[(iii)] If \( \{ {C}_{i} \}_{i \in I }  \) is a collection of closed sets, then \( \bigcap_{  i \in I  }^{  }  {C}_{i} \) is closed.
\end{enumerate}
\begin{remark}
    Condition (ii) is equivalent to ii'): if \( {C}_{1}, {C}_{2}, \dots, {C}_{n} \) are closed, then \( {C}_{1} \cup {C}_{2} \cup \cdots \cup {C}_{n} \) is closed.
\end{remark}
\end{prop}
\begin{proof}
\textit{left to the reader.} 
\end{proof}

\begin{definition}[Points of Closure (adherent point)]
    \begin{itemize}
        \item We say that \( x  \) is a \textbf{point of closure} (or an \textbf{adherent point}) of \( A  \) if
            \[  \forall N \in \mathcal{N}(x), \ \text{we have} \ N \cap A \neq \emptyset.  \]
        \item The \textbf{closure} of \( A  \) is defined as
            \[  \overline{A} = \{ {x}_{0} \in X : {x}_{0} \ \text{is a closure point of \( A  \)} \}. \]
    \end{itemize}
\end{definition}

\begin{remark}
    We can express these ideas in terms of neighborhoods to open sets and to basis elements (assuming that a basis exists for the topology).
\end{remark}

Let's prove the third observation.

\begin{proof}
\( (\Longrightarrow) \) Let \( x \in \overline{A} \). Let \( b \in \mathcal{B} \) such that \( x \in B \). Recall that \( B \) is open (by definition of a basis), of \( B \in \mathcal{N}(x) \). By definition of point of closure, we get that \( B \cap A \neq \emptyset  \).

\( (\Longleftarrow) \) 
\end{proof}

\begin{remark} Let \( (X,\mathcal{T}) \) be a topological space and \( A \subseteq X  \). Then 
   \begin{itemize}
       \item \( A \subseteq \overline{A} \).
        \item \( x \in \overline{A }  \) if and only if \( U \cap A \neq \emptyset  \) for every open set \( U  \) such that \( x \in U  \). 
        \item If \( \mathcal{B} \) is a basis for \( \mathcal{T} \), then \( x \in \overline{A}\) if and only if \( B \cap A \neq \emptyset  \) for all \( B \in \mathcal{B} \) such that \( x \in B \).
   \end{itemize} 
\end{remark}

\begin{prop}
    Let \( X  \) be a topological space and \( A \subseteq X  \). We have that 
    \begin{center}
        \( A  \) is closed if and only if \( A = \overline{A} \).
    \end{center}
\end{prop}

\begin{prop}
    Let \( X  \) be a topological space and \( A \subseteq X  \). The closure of \( A  \) is the smallest closed set containing \( A  \). Equivalently, 
    \[  \overline{A} = \bigcap_{ C \in \Omega }^{  } C \]
    where \( \Omega = \{  A \subseteq C : \text{\( C \) is a closed set} \}  \).
\end{prop}

\begin{problem}
    For \( a < b \in \R  \) consider \( A = (a,b) \), \( B = [a,b) \), \( C = [a,b] \), and \( D = \Q  \). Find the closure of \( A , B , C , D  \) as subsets of \( \R  \) endowed with
    \begin{enumerate}
        \item[(i)] The Euclidean Topology,
        \item[(ii)] The lower limit topology,
        \item[(iii)] The cofinite topology.
    \end{enumerate}
\end{problem}

\begin{definition}[Dense]
    Let \( X  \) be a topological space and \( A \subseteq X   \). We say that \( A  \) is \textbf{dense} in \( X  \) if \( \overline{A} = X  \).
\end{definition}

\begin{eg}
    \begin{enumerate}
        \item[(i)] The set of rationals \( \Q  \) in \( \R  \) equipped with the Euclidean topology (and with the lower limit topology).
        \item[(ii)] More in general, for \( n \in \N  \), the set \( \Q^{n}  \) is dense in \( \R^{n} \) equipped with the Euclidean topology.
        \item[(iii)] The set of irrationals \( \R \setminus  \Q  \) is dense in \( \R  \) equipped with the Euclidean topology.
    \end{enumerate}
\end{eg}

\begin{prop}
    Let \( (X,\mathcal{T}) \) be a topological space and \( A \subseteq X  \). We have that 
    \begin{enumerate}
        \item[(i)] \( A  \) is dense in \( X  \) if and only if \( A \cap U \neq \emptyset \) for all nonempty \( U \in \mathcal{T} \);
        \item[(ii)] If \( \mathcal{B} \) is basis for \( \mathcal{T} \), \( A  \) is dense in \( X  \), if and only if \( A \cap B \neq \emptyset \) for all nonempty \( B \in \mathcal{B} \).
    \end{enumerate}
\end{prop}
\begin{proof}
    \begin{enumerate}
        \item[(i)] Assume that \( A  \) is dense in \( X  \); that is, \( \overline{A} = X  \). Let \( \emptyset \neq U \in \mathcal{T} \). Then there exists \( x \in U  \). By assumption, \( x \in \overline{A} \). Since \( U  \) is open and \( x \in U  \), then \( U \in \mathcal{N}(x) \) so \( U \cap A \neq \emptyset \), by definition of closure.
        Conversely, assume that \( A \cap U \neq \emptyset \) for all \( U \in \mathcal{T} \), \( U \neq \emptyset \). Let \( x \in X  \). Consider an open set \( U  \) containing \( x  \). By assumption, \( U \cap A \neq \emptyset \), so \( x \in \overline{A} \). This shows that \( X \subseteq \overline{A} \). Since \( \overline{A} \subseteq X  \), we conclude that \( \overline{A} = X  \) and so \( A  \) must be dense in \( X  \).
        \item[(ii)] \textit{Left to reader as an exercise.} 
    \end{enumerate}
\end{proof}


